\section{PCA Pre-processing and Conclusion}

\subsection*{Use PCA to reduce the dimensionality of the feature space, and apply the three classification approaches. What do you observe?}
Based on the experimental results, we observe a steady increase in the error rate as the number of PCA dimensions ($m$) is reduced. For example, the Standard MVG error rate increases from its baseline to 8.8\% at $m=2$, and further to 9.3\% at $m=1$. This indicates that reducing the subspace dimensionality leads to a loss of information that is critical for classification.

\subsection*{Is PCA effective for this dataset with the Gaussian models?}
No, PCA is not effective as a pre-processing step for this specific dataset. PCA discards dimensions with lower variance; however, the results show that these discarded dimensions contain significant discriminative information. Consequently, any reduction in dimensionality through PCA results in a higher error rate compared to using the original 6-dimensional feature space.

\subsection*{Overall, which is the model that provided the best accuracy on the validation set?}
The model providing the highest accuracy (and lowest error rate) on the validation set is the \textbf{Standard Multivariate Gaussian (MVG)} using the \textbf{full 6-dimensional feature set}. This model achieved an optimal error rate of \textbf{7.0\%}, outperforming both the Tied Covariance and Naive Bayes variations, as well as all PCA-reduced versions of the models.